广泛扫描通过非对称距离计算评估主码重建：查询保持 FP32，存储编码用于选择重建坐标\cite{pq,jhq}。对变换后的查询 $q'$，主码子空间 $m$ 的普通查找表为
\begin{equation}
 T_m[z]=\|q'_m-C_m[z]\|_2^2,\qquad
 \primarydist(q,x)=\sum_{m=0}^{M-1}T_m[z_m(x)].
 \label{eq:5-primary-lut}
\end{equation}
八比特码要求每个子空间存储 256 项。这里的优化机会来自码本内部的 Cartesian 结构，而不只是向量被划分成 PQ 子空间。

\subsection{在不改变量化器的情况下分解}

对被测 $D_s=8$ Cartesian 码本，高半字节 $h$ 选择前四个重建坐标，低半字节 $l$ 选择后四个。因此 $C_m[16h+l]=(A_m[h],B_m[l])$，两个坐标组之间不存在交叉依赖。定义
\begin{equation}
 T_m^{\rm hi}[h]=\|q'_{m,\rm hi}-A_m[h]\|_2^2,\qquad
 T_m^{\rm lo}[l]=\|q'_{m,\rm lo}-B_m[l]\|_2^2.
 \label{eq:5-lut-halves}
\end{equation}
平方 L2 对互不相交坐标可加，得到
\begin{equation}
 T_m[z]=T_m^{\rm hi}[z\!\gg\!4]
        +T_m^{\rm lo}[z\mathbin{\&}15].
 \label{eq:5-lut-factorisation}
\end{equation}
只需存储 32 项，而不是 256 项。重建表示和实数运算下的分数保持不变，无需重训练编码、增加量化或改变候选集合。

建表程序从现有码本读取代表码字 $16h$ 和 $l$，只计算相应坐标片段；另一片段的零编码位与此计算无关。这要求所有共享同一半字节的码字也共享对应坐标片段。\texttt{verify\_split\_lut} 将码字片段与代表码字比较，拒绝不兼容码本。任意学习得到的 256 中心 PQ 子码本不一定满足该条件；仅将编码标签拆成半字节，不会自动分解中心坐标。

数学恒等式不保证 FP32 逐位相等。八个坐标的连续累加，与两个四坐标和再相加，浮点结合顺序不同。因此实现保留 FP32 表，并单独检查数值和排序影响。将表项改成 FP16 会增加另一种舍入，是不同的设计选择。

\subsection{分组、建表工作量与资源成本}

v71 消融将分解推广到 $g=\glut\in\{1,2,4,8\}$ 个连续坐标组，每组位字段与 Cartesian 数字对应。当编码为八比特、$D_s=8$ 时，每组有 $2^{8/g}$ 项；总表项和每查询 FP32 存储为
\begin{equation}
 E(g)=g2^{8/g},\qquad W_{\rm LUT}(g)=4ME(g)\ \text{bytes}.
 \label{eq:5-lut-cost}
\end{equation}
独立构造表项时，每个子空间计算 $D_s2^{8/g}$ 个坐标贡献。相比独立计算完整表的每一项，两组设计将表存储减少八倍、坐标工作量减少十六倍。这是表级成本比，不是应用性能加速预测。

\begin{table}[tbp]
\centering
\caption{八比特、八坐标主码子空间的 Cartesian 表权衡。合并加法为组合查表结果所需的代数最小次数，不含子空间之间的累加，也不是实测指令数。}
\label{tab:5-lut-resources}
\small
\begin{tabular}{@{}rrrrrr@{}}
\toprule
$\glut$ & 表项 & FP32 字节 & 查表次数 & 合并加法 & 建表坐标项 \\
\midrule
1 & 256 & 1,024 & 1 & 0 & 2,048 \\
2 & 32 & 128 & 2 & 1 & 128 \\
4 & 16 & 64 & 4 & 3 & 32 \\
8 & 16 & 64 & 8 & 7 & 16 \\
\bottomrule
\end{tabular}
\end{table}

表先在全局工作区构造。如果表和该 kernel 的候选缓冲区能放入设备配置的容量，扫描 block 就合作将表搬到共享内存，再通过 block 屏障完成搬运；否则读取全局表。共享内存占用必须包含选择工作区，不能仅由 $W_{\rm LUT}$ 判断驻留。多个分组扫描 block 还可能各自搬运同一查询的表，因此这一设计与\Cref{sec:5-5}的选择有关。

当 $M=96$ 和 384 时，每查询完整 FP32 表分别需要 96 和 384\,KiB，两组表仅需 12 和 48\,KiB。较小占用提供更多放置机会，但四组和八组即使均只存储每子空间 16 项，也要做更多查表和结果组合。索引提取、依赖查表、同步和驻留均参与最终权衡。不能仅凭表大小判断最快设计，或确定唯一的 cache／bank-conflict 原因。

小表执行在 PQ Fast Scan 和 Quicker ADC 中已有重要先例\cite{fastscan,quickeradc}。这里将精确坐标分解应用于 JHQ 现有 Cartesian 码本，并测量 GPU 分组权衡，不宣称首次提出小 LUT。\Cref{sec:7-3}报告受控 v71 扫描：两组在 24 个单元中领先 22 个，另外两组差距均不足 1\%。结果支持选定运行点，而非普遍最优结论。
